\section{}

\paragraph{1057年}\eventanchor{NS-1057-EXAMINATION}{NS-1057-EXAMINATION}　礼部进士试由欧阳修等主持，取士文风和士人网络发生显著变化。账本所列《宋史》卷11〈仁宗本纪四〉；《宋史》卷319〈欧阳修传〉；《宋会要辑稿·选举》嘉祐二年进士条为本条来源起点；该条可固定编年位置，具体执行、规模和地方社会影响仍须保留来源差异。

\paragraph{1057年}\eventanchor{NS-1057-REGNAL-YEAR}{NS-1057-REGNAL-YEAR}　宋以宋仁宗在位第35年作为诏令、赋役与外交文书的纪年框架。账本所列《宋史》卷1—23〈本纪〉；《续资治通鉴长编》本年条；《宋会要辑稿》相关门类为本条来源起点；该条可固定编年位置，具体执行、规模和地方社会影响仍须保留来源差异。

\paragraph{1058年}\eventanchor{NS-1058-REGNAL-YEAR}{NS-1058-REGNAL-YEAR}　宋以宋仁宗在位第36年作为诏令、赋役与外交文书的纪年框架。账本所列《宋史》卷1—23〈本纪〉；《续资治通鉴长编》本年条；《宋会要辑稿》相关门类为本条来源起点；该条可固定编年位置，具体执行、规模和地方社会影响仍须保留来源差异。

\paragraph{1059年}\eventanchor{NS-1059-REGNAL-YEAR}{NS-1059-REGNAL-YEAR}　宋以宋仁宗在位第37年作为诏令、赋役与外交文书的纪年框架。账本所列《宋史》卷1—23〈本纪〉；《续资治通鉴长编》本年条；《宋会要辑稿》相关门类为本条来源起点；该条可固定编年位置，具体执行、规模和地方社会影响仍须保留来源差异。

\paragraph{1060年}\eventanchor{NS-1060-REGNAL-YEAR}{NS-1060-REGNAL-YEAR}　宋以宋仁宗在位第38年作为诏令、赋役与外交文书的纪年框架。账本所列《宋史》卷1—23〈本纪〉；《续资治通鉴长编》本年条；《宋会要辑稿》相关门类为本条来源起点；该条可固定编年位置，具体执行、规模和地方社会影响仍须保留来源差异。

\paragraph{1061年}\eventanchor{NS-1061-REGNAL-YEAR}{NS-1061-REGNAL-YEAR}　宋以宋仁宗在位第39年作为诏令、赋役与外交文书的纪年框架。账本所列《宋史》卷1—23〈本纪〉；《续资治通鉴长编》本年条；《宋会要辑稿》相关门类为本条来源起点；该条可固定编年位置，具体执行、规模和地方社会影响仍须保留来源差异。

\paragraph{1062年}\eventanchor{NS-1062-REGNAL-YEAR}{NS-1062-REGNAL-YEAR}　宋以宋仁宗在位第40年作为诏令、赋役与外交文书的纪年框架。账本所列《宋史》卷1—23〈本纪〉；《续资治通鉴长编》本年条；《宋会要辑稿》相关门类为本条来源起点；该条可固定编年位置，具体执行、规模和地方社会影响仍须保留来源差异。

\paragraph{1063年}\eventanchor{NS-1063-EVENT-003}{NS-1063-EVENT-003}　仁宗去世，英宗即位。账本所列《宋史》卷1—23〈本纪〉；《续资治通鉴长编》本年条；《宋会要辑稿》相关门类为本条来源起点；该条可固定编年位置，具体执行、规模和地方社会影响仍须保留来源差异。

\paragraph{1063年}\eventanchor{NS-1063-REGNAL-YEAR}{NS-1063-REGNAL-YEAR}　宋以宋仁宗在位第41年作为诏令、赋役与外交文书的纪年框架。账本所列《宋史》卷1—23〈本纪〉；《续资治通鉴长编》本年条；《宋会要辑稿》相关门类为本条来源起点；该条可固定编年位置，具体执行、规模和地方社会影响仍须保留来源差异。

\paragraph{1063年三月}\eventanchor{NS-1063-RENZONG-DEATH}{NS-1063-RENZONG-DEATH}　仁宗去世，英宗继位。账本所列《宋史》卷12〈仁宗本纪五〉；《宋史》卷13〈英宗本纪〉；《续资治通鉴长编》嘉祐八年三月条为本条来源起点；该条可固定编年位置，具体执行、规模和地方社会影响仍须保留来源差异。

\paragraph{1064年}\eventanchor{NS-1064-EVENT-004}{NS-1064-EVENT-004}　围绕英宗生父濮王名分的濮议在朝廷展开。账本所列《宋史》卷1—23〈本纪〉；《续资治通鉴长编》本年条；《宋会要辑稿》相关门类为本条来源起点；该条可固定编年位置，具体执行、规模和地方社会影响仍须保留来源差异。

\paragraph{1064年}\eventanchor{NS-1064-REGNAL-YEAR}{NS-1064-REGNAL-YEAR}　宋以宋英宗在位第1年作为诏令、赋役与外交文书的纪年框架。账本所列《宋史》卷1—23〈本纪〉；《续资治通鉴长编》本年条；《宋会要辑稿》相关门类为本条来源起点；该条可固定编年位置，具体执行、规模和地方社会影响仍须保留来源差异。

\paragraph{1065年}\eventanchor{NS-1065-PUYI-DEBATE}{NS-1065-PUYI-DEBATE}　围绕英宗生父濮王的追尊发生濮议。账本所列《宋史》卷13〈英宗本纪〉；《宋史》卷336〈司马光传〉；《续资治通鉴长编》治平二年条为本条来源起点；该条可固定编年位置，具体执行、规模和地方社会影响仍须保留来源差异。

\paragraph{1065年}\eventanchor{NS-1065-REGNAL-YEAR}{NS-1065-REGNAL-YEAR}　宋以宋英宗在位第2年作为诏令、赋役与外交文书的纪年框架。账本所列《宋史》卷1—23〈本纪〉；《续资治通鉴长编》本年条；《宋会要辑稿》相关门类为本条来源起点；该条可固定编年位置，具体执行、规模和地方社会影响仍须保留来源差异。

\paragraph{1066年}\eventanchor{NS-1066-REGNAL-YEAR}{NS-1066-REGNAL-YEAR}　宋以宋英宗在位第3年作为诏令、赋役与外交文书的纪年框架。账本所列《宋史》卷1—23〈本纪〉；《续资治通鉴长编》本年条；《宋会要辑稿》相关门类为本条来源起点；该条可固定编年位置，具体执行、规模和地方社会影响仍须保留来源差异。

\paragraph{1067年}\eventanchor{NS-1067-EVENT-005}{NS-1067-EVENT-005}　英宗去世，神宗即位。账本所列《宋史》卷1—23〈本纪〉；《续资治通鉴长编》本年条；《宋会要辑稿》相关门类为本条来源起点；该条可固定编年位置，具体执行、规模和地方社会影响仍须保留来源差异。

\paragraph{1067年}\eventanchor{NS-1067-REGNAL-YEAR}{NS-1067-REGNAL-YEAR}　宋以宋英宗在位第4年作为诏令、赋役与外交文书的纪年框架。账本所列《宋史》卷1—23〈本纪〉；《续资治通鉴长编》本年条；《宋会要辑稿》相关门类为本条来源起点；该条可固定编年位置，具体执行、规模和地方社会影响仍须保留来源差异。

\paragraph{1067年正月}\eventanchor{NS-1067-SHENZONG-ACCESSION}{NS-1067-SHENZONG-ACCESSION}　英宗去世，神宗继位。账本所列《宋史》卷13〈英宗本纪〉；《宋史》卷14〈神宗本纪一〉；《续资治通鉴长编》治平四年正月条为本条来源起点；该条可固定编年位置，具体执行、规模和地方社会影响仍须保留来源差异。

\paragraph{1068年}\eventanchor{NS-1068-EVENT-006}{NS-1068-EVENT-006}　王安石受神宗倚重进入中枢。账本所列《宋史》卷1—23〈本纪〉；《续资治通鉴长编》本年条；《宋会要辑稿》相关门类为本条来源起点；该条可固定编年位置，具体执行、规模和地方社会影响仍须保留来源差异。

\paragraph{1068年}\eventanchor{NS-1068-REGNAL-YEAR}{NS-1068-REGNAL-YEAR}　宋以宋神宗在位第1年作为诏令、赋役与外交文书的纪年框架。账本所列《宋史》卷1—23〈本纪〉；《续资治通鉴长编》本年条；《宋会要辑稿》相关门类为本条来源起点；该条可固定编年位置，具体执行、规模和地方社会影响仍须保留来源差异。

\paragraph{1068年}\eventanchor{NS-1068-WANG-ANSHI-COURT}{NS-1068-WANG-ANSHI-COURT}　神宗召用王安石并与其讨论变法方向。账本所列《宋史》卷14〈神宗本纪一〉；《宋史》卷327〈王安石传〉；《续资治通鉴长编》熙宁元年条为本条来源起点；该条可固定编年位置，具体执行、规模和地方社会影响仍须保留来源差异。

\paragraph{1069年}\eventanchor{NS-1069-EVENT-007}{NS-1069-EVENT-007}　熙宁新法陆续推行，青苗、均输等成为改革核心。账本所列《宋史》卷1—23〈本纪〉；《续资治通鉴长编》本年条；《宋会要辑稿》相关门类为本条来源起点；该条可固定编年位置，具体执行、规模和地方社会影响仍须保留来源差异。

\paragraph{1069年}\eventanchor{NS-1069-GREEN-SPROUTS}{NS-1069-GREEN-SPROUTS}　青苗法颁行，以官贷介入青黄不接时的借贷。账本所列《宋史》卷175〈食货志上三〉；《宋会要辑稿·食货》熙宁二年青苗条；《续资治通鉴长编》熙宁二年条为本条来源起点；该条可固定编年位置，具体执行、规模和地方社会影响仍须保留来源差异。

\paragraph{1069年}\eventanchor{NS-1069-REGNAL-YEAR}{NS-1069-REGNAL-YEAR}　宋以宋神宗在位第2年作为诏令、赋役与外交文书的纪年框架。账本所列《宋史》卷1—23〈本纪〉；《续资治通鉴长编》本年条；《宋会要辑稿》相关门类为本条来源起点；该条可固定编年位置，具体执行、规模和地方社会影响仍须保留来源差异。

\paragraph{1070年}\eventanchor{NS-1070-EXEMPTION-SERVICE}{NS-1070-EXEMPTION-SERVICE}　免役法在部分地区推进，以雇募和免役钱调整差役。账本所列《宋史》卷178〈食货志下一〉；《宋会要辑稿·食货》熙宁免役条；《续资治通鉴长编》熙宁三年条为本条来源起点；该条可固定编年位置，具体执行、规模和地方社会影响仍须保留来源差异。

\paragraph{1070年}\eventanchor{NS-1070-REGNAL-YEAR}{NS-1070-REGNAL-YEAR}　宋以宋神宗在位第3年作为诏令、赋役与外交文书的纪年框架。账本所列《宋史》卷1—23〈本纪〉；《续资治通鉴长编》本年条；《宋会要辑稿》相关门类为本条来源起点；该条可固定编年位置，具体执行、规模和地方社会影响仍须保留来源差异。

\paragraph{1071年}\eventanchor{NS-1071-BAOJIA}{NS-1071-BAOJIA}　保甲法开始推行，以户籍编组承担治安、训练和互保功能。账本所列《宋史》卷192〈兵志六〉；《宋会要辑稿·兵》熙宁保甲条；《续资治通鉴长编》熙宁四年条为本条来源起点；该条可固定编年位置，具体执行、规模和地方社会影响仍须保留来源差异。

\paragraph{1071年}\eventanchor{NS-1071-EVENT-008}{NS-1071-EVENT-008}　免役法在部分地区推行，以募役钱调整差役结构。账本所列《宋史》卷1—23〈本纪〉；《续资治通鉴长编》本年条；《宋会要辑稿》相关门类为本条来源起点；该条可固定编年位置，具体执行、规模和地方社会影响仍须保留来源差异。

\paragraph{1071年}\eventanchor{NS-1071-REGNAL-YEAR}{NS-1071-REGNAL-YEAR}　宋以宋神宗在位第4年作为诏令、赋役与外交文书的纪年框架。账本所列《宋史》卷1—23〈本纪〉；《续资治通鉴长编》本年条；《宋会要辑稿》相关门类为本条来源起点；该条可固定编年位置，具体执行、规模和地方社会影响仍须保留来源差异。

\paragraph{1072年}\eventanchor{NS-1072-MARKET-EXCHANGE}{NS-1072-MARKET-EXCHANGE}　市易法设市易务，以官本调节城市商货和信贷。账本所列《宋史》卷181〈食货志下三〉；《宋会要辑稿·食货》熙宁市易条；《续资治通鉴长编》熙宁五年条为本条来源起点；该条可固定编年位置，具体执行、规模和地方社会影响仍须保留来源差异。

\paragraph{1072年}\eventanchor{NS-1072-REGNAL-YEAR}{NS-1072-REGNAL-YEAR}　宋以宋神宗在位第5年作为诏令、赋役与外交文书的纪年框架。账本所列《宋史》卷1—23〈本纪〉；《续资治通鉴长编》本年条；《宋会要辑稿》相关门类为本条来源起点；该条可固定编年位置，具体执行、规模和地方社会影响仍须保留来源差异。

\paragraph{1073年}\eventanchor{NS-1073-REGNAL-YEAR}{NS-1073-REGNAL-YEAR}　宋以宋神宗在位第6年作为诏令、赋役与外交文书的纪年框架。账本所列《宋史》卷1—23〈本纪〉；《续资治通鉴长编》本年条；《宋会要辑稿》相关门类为本条来源起点；该条可固定编年位置，具体执行、规模和地方社会影响仍须保留来源差异。

\paragraph{1074年}\eventanchor{NS-1074-DROUGHT-ANSHI-LEAVES}{NS-1074-DROUGHT-ANSHI-LEAVES}　旱灾、政治批评与人事冲突交织，王安石首次罢相离京。账本所列《宋史》卷15〈神宗本纪二〉；《宋史》卷327〈王安石传〉；《续资治通鉴长编》熙宁七年条为本条来源起点；该条可固定编年位置，具体执行、规模和地方社会影响仍须保留来源差异。

\paragraph{1074年}\eventanchor{NS-1074-EVENT-009}{NS-1074-EVENT-009}　灾荒与政论交织，反对者以民间困乏批评新法执行。账本所列《宋史》卷1—23〈本纪〉；《续资治通鉴长编》本年条；《宋会要辑稿》相关门类为本条来源起点；该条可固定编年位置，具体执行、规模和地方社会影响仍须保留来源差异。

\paragraph{1074年}\eventanchor{NS-1074-REGNAL-YEAR}{NS-1074-REGNAL-YEAR}　宋以宋神宗在位第7年作为诏令、赋役与外交文书的纪年框架。账本所列《宋史》卷1—23〈本纪〉；《续资治通鉴长编》本年条；《宋会要辑稿》相关门类为本条来源起点；该条可固定编年位置，具体执行、规模和地方社会影响仍须保留来源差异。

\paragraph{1075年}\eventanchor{NS-1075-EVENT-010}{NS-1075-EVENT-010}　宋军在熙河方向推进，经营堡寨、道路与羁縻关系。账本所列《宋史》卷1—23〈本纪〉；《续资治通鉴长编》本年条；《宋会要辑稿》相关门类为本条来源起点；该条可固定编年位置，具体执行、规模和地方社会影响仍须保留来源差异。

\paragraph{1075年冬}\eventanchor{NS-1075-LY-INVASION}{NS-1075-LY-INVASION}　李朝军队攻入邕州、钦州、廉州等地。账本所列《宋史》卷15〈神宗本纪二〉；《宋史》卷488〈交阯传〉；《大越史记全书》李朝条为本条来源起点；该条可固定编年位置，具体执行、规模和地方社会影响仍须保留来源差异。

\paragraph{1075年}\eventanchor{NS-1075-REGNAL-YEAR}{NS-1075-REGNAL-YEAR}　宋以宋神宗在位第8年作为诏令、赋役与外交文书的纪年框架。账本所列《宋史》卷1—23〈本纪〉；《续资治通鉴长编》本年条；《宋会要辑稿》相关门类为本条来源起点；该条可固定编年位置，具体执行、规模和地方社会影响仍须保留来源差异。

\paragraph{1076—1077年}\eventanchor{NS-1076-GUANGNAN-COUNTEROFFENSIVE}{NS-1076-GUANGNAN-COUNTEROFFENSIVE}　郭逵等率军南征，战至富良江一线后撤。账本所列《宋史》卷15〈神宗本纪二〉；《宋史》卷488〈交阯传〉；《续资治通鉴长编》熙宁九年至十年条为本条来源起点；该条可固定编年位置，具体执行、规模和地方社会影响仍须保留来源差异。

\paragraph{1076年}\eventanchor{NS-1076-REGNAL-YEAR}{NS-1076-REGNAL-YEAR}　宋以宋神宗在位第9年作为诏令、赋役与外交文书的纪年框架。账本所列《宋史》卷1—23〈本纪〉；《续资治通鉴长编》本年条；《宋会要辑稿》相关门类为本条来源起点；该条可固定编年位置，具体执行、规模和地方社会影响仍须保留来源差异。

\paragraph{1077年}\eventanchor{NS-1077-REGNAL-YEAR}{NS-1077-REGNAL-YEAR}　宋以宋神宗在位第10年作为诏令、赋役与外交文书的纪年框架。账本所列《宋史》卷1—23〈本纪〉；《续资治通鉴长编》本年条；《宋会要辑稿》相关门类为本条来源起点；该条可固定编年位置，具体执行、规模和地方社会影响仍须保留来源差异。

\paragraph{1078年}\eventanchor{NS-1078-REGNAL-YEAR}{NS-1078-REGNAL-YEAR}　宋以宋神宗在位第11年作为诏令、赋役与外交文书的纪年框架。账本所列《宋史》卷1—23〈本纪〉；《续资治通鉴长编》本年条；《宋会要辑稿》相关门类为本条来源起点；该条可固定编年位置，具体执行、规模和地方社会影响仍须保留来源差异。

\paragraph{1079年}\eventanchor{NS-1079-EVENT-011}{NS-1079-EVENT-011}　苏轼因诗文案受审，后出知地方。账本所列《宋史》卷1—23〈本纪〉；《续资治通鉴长编》本年条；《宋会要辑稿》相关门类为本条来源起点；该条可固定编年位置，具体执行、规模和地方社会影响仍须保留来源差异。

\paragraph{1079年}\eventanchor{NS-1079-REGNAL-YEAR}{NS-1079-REGNAL-YEAR}　宋以宋神宗在位第12年作为诏令、赋役与外交文书的纪年框架。账本所列《宋史》卷1—23〈本纪〉；《续资治通鉴长编》本年条；《宋会要辑稿》相关门类为本条来源起点；该条可固定编年位置，具体执行、规模和地方社会影响仍须保留来源差异。

\paragraph{1079年}\eventanchor{NS-1079-WUTAI-POETRY}{NS-1079-WUTAI-POETRY}　苏轼因诗文被控讥讽朝政，发生乌台诗案。账本所列《宋史》卷338〈苏轼传〉；《续资治通鉴长编》元丰二年条；《宋会要辑稿·刑法》言事条为本条来源起点；该条可固定编年位置，具体执行、规模和地方社会影响仍须保留来源差异。

\paragraph{1080年}\eventanchor{NS-1080-EVENT-012}{NS-1080-EVENT-012}　永乐城陷落，宋军在西北遭受重大挫败。账本所列《宋史》卷1—23〈本纪〉；《续资治通鉴长编》本年条；《宋会要辑稿》相关门类为本条来源起点；该条可固定编年位置，具体执行、规模和地方社会影响仍须保留来源差异。

\paragraph{1080年}\eventanchor{NS-1080-REGNAL-YEAR}{NS-1080-REGNAL-YEAR}　宋以宋神宗在位第13年作为诏令、赋役与外交文书的纪年框架。账本所列《宋史》卷1—23〈本纪〉；《续资治通鉴长编》本年条；《宋会要辑稿》相关门类为本条来源起点；该条可固定编年位置，具体执行、规模和地方社会影响仍须保留来源差异。

\paragraph{1080年}\eventanchor{NS-1080-YONGLE}{NS-1080-YONGLE}　宋军筑永乐城并遭西夏围攻失守。账本所列《宋史》卷15〈神宗本纪二〉；《宋史》卷485〈夏国传上〉；《续资治通鉴长编》元丰三年条为本条来源起点；该条可固定编年位置，具体执行、规模和地方社会影响仍须保留来源差异。

\paragraph{1081年}\eventanchor{NS-1081-EVENT-013}{NS-1081-EVENT-013}　宋以多路军攻夏，战事未取得预期结果。账本所列《宋史》卷1—23〈本纪〉；《续资治通鉴长编》本年条；《宋会要辑稿》相关门类为本条来源起点；该条可固定编年位置，具体执行、规模和地方社会影响仍须保留来源差异。

\paragraph{1081年}\eventanchor{NS-1081-REGNAL-YEAR}{NS-1081-REGNAL-YEAR}　宋以宋神宗在位第14年作为诏令、赋役与外交文书的纪年框架。账本所列《宋史》卷1—23〈本纪〉；《续资治通鉴长编》本年条；《宋会要辑稿》相关门类为本条来源起点；该条可固定编年位置，具体执行、规模和地方社会影响仍须保留来源差异。
